\section{Controlled generator and paired benchmark}
\label{sec:benchmark}

\subsection{Renderer contract}

% receipt: evidence/benchmark_v2/design.plan.json; source_sha256 renderer and generator entries
% source: rendering/generate_benchmark_v2.py; rendering/benchmark_v2.js
Each scene is rendered at $1920\times1080$ pixels from a procedural board, a
piece family with a recorded asset grant, a perspective camera, and explicit lighting and material
parameters. The manifest stores the position in Forsyth--Edwards Notation, the
camera state, the factor assignment, and the projected board corners. Labels
therefore come from the scene description rather than from a detector or a
recognizer. The labels describe the latent scene state; they do not certify that
every occupied square is visually identifiable in a single image under occlusion.

The renderer treats board colour convention as an explicit scene variable.
% receipt: evidence/benchmark_v2/benchmark.summary.json
The reference condition makes a1 dark; its paired stress condition makes a1
light while retaining the position, camera, geometry, palette, and framing.
This intervention changes the physical checkerboard colours, not the label
order or the image orientation.

\subsection{Paired factor design}

% receipt: evidence/benchmark_v2/benchmark.summary.json
The benchmark varies four elevation bands, eight piece geometries, six non-reference
material schemes, two framing ranges, and two board conventions. These settings form
factor panels rather than a full Cartesian product: each non-reference level is
paired with the appropriate reference scene, and every scene has both convention
mates. A palette treatment is a supplied material scheme: it changes board and
piece colours and may also change roughness and metalness. It is not an RGB-only
intervention. The reference uses the highest elevation band, Poly Haven geometry,
standard framing, and the separate reference material scheme. Table~\ref{tab:design} gives the registered levels; Figure~\ref{fig:design}
summarizes the data path and the pairing structure. These are contrasts conditional
on the reference settings; the design does not estimate all factor interactions.

\begingroup\small
\setlength{\tabcolsep}{4pt}
\begin{table}[htbp]\centering
\caption{Paired factor panels, not a full factorial cross-product.}\label{tab:design}
\begin{tabular}{>{\raggedright\arraybackslash}p{.16\linewidth}r>{\raggedright\arraybackslash}p{.65\linewidth}}
\toprule
Factor & % receipt: evidence/benchmark_v2/benchmark.complete.json; source: evidence/benchmark_v2/benchmark.summary.json; column definitions
Levels & % receipt: evidence/benchmark_v2/benchmark.complete.json; source: evidence/benchmark_v2/benchmark.summary.json; column definitions
Values \\ % receipt: evidence/benchmark_v2/benchmark.complete.json; source: evidence/benchmark_v2/benchmark.summary.json; column definitions
\midrule

Elevation & % receipt: evidence/benchmark_v2/benchmark.complete.json; source: evidence/benchmark_v2/benchmark.summary.json; cells/elevation
4 & % receipt: evidence/benchmark_v2/benchmark.complete.json; source: evidence/benchmark_v2/benchmark.summary.json; cells/elevation
08\_\allowbreak{}15, 15\_\allowbreak{}25, 25\_\allowbreak{}45, 45\_\allowbreak{}60 \\ % receipt: evidence/benchmark_v2/benchmark.complete.json; source: evidence/benchmark_v2/benchmark.summary.json; cells/elevation

Geometry & % receipt: evidence/benchmark_v2/benchmark.complete.json; source: evidence/benchmark_v2/benchmark.summary.json; cells/geometry
8 & % receipt: evidence/benchmark_v2/benchmark.complete.json; source: evidence/benchmark_v2/benchmark.summary.json; cells/geometry
deplorable, killgorack, lyricsz, nailfighter, omie, phuong, pjhanzlik, polyhaven \\ % receipt: evidence/benchmark_v2/benchmark.complete.json; source: evidence/benchmark_v2/benchmark.summary.json; cells/geometry

Palette & % receipt: evidence/benchmark_v2/benchmark.complete.json; source: evidence/benchmark_v2/benchmark.summary.json; cells/palette
6 & % receipt: evidence/benchmark_v2/benchmark.complete.json; source: evidence/benchmark_v2/benchmark.summary.json; cells/palette
aubergine\_\allowbreak{}periwinkle, cobalt\_\allowbreak{}porcelain, indigo\_\allowbreak{}lavender, navy\_\allowbreak{}coolrose, plum\_\allowbreak{}mintgreen, wine\_\allowbreak{}iceblue \\ % receipt: evidence/benchmark_v2/benchmark.complete.json; source: evidence/benchmark_v2/benchmark.summary.json; cells/palette

Framing & % receipt: evidence/benchmark_v2/benchmark.complete.json; source: evidence/benchmark_v2/benchmark.summary.json; cells/framing
2 & % receipt: evidence/benchmark_v2/benchmark.complete.json; source: evidence/benchmark_v2/benchmark.summary.json; cells/framing
compact, standard \\ % receipt: evidence/benchmark_v2/benchmark.complete.json; source: evidence/benchmark_v2/benchmark.summary.json; cells/framing

Convention & % receipt: evidence/benchmark_v2/benchmark.complete.json; source: evidence/benchmark_v2/benchmark.summary.json; cells/parity
2 & % receipt: evidence/benchmark_v2/benchmark.complete.json; source: evidence/benchmark_v2/benchmark.summary.json; cells/parity
a1 dark, a1 light \\ % receipt: evidence/benchmark_v2/benchmark.complete.json; source: evidence/benchmark_v2/benchmark.summary.json; cells/parity

\bottomrule
\end{tabular}
\end{table}
\endgroup


% receipt: evidence/benchmark_v2/benchmark.summary.json
The design contains 500 independently sampled position groups. A fixed split
assigns 400 groups to development and 100 to the sealed final test, with every
view of a group kept in the same partition. Applying the 20 factor cells at
both conventions yields 20,000 cell-images: 16,000 development images and
4,000 sealed images. Three reference cells deliberately reuse identical scene
configurations, so the benchmark contains 18,000 unique PNG hashes. Neither
cell-images nor repeated render hashes are independent board samples; the
independent sampling unit remains the position group.

% receipt: evidence/benchmark_v2/design.plan.json; source_sha256
% source: rendering/generate_benchmark_v2.py:generate_groups; rendering/generate_geometry.py:play_position
Each group contributes one position from a separately simulated legal random-play
game. The generator deliberately mixes dense and sparse positions and biases move
selection toward captures. It is not a sample of human game frequencies. Exact
and symmetry checks do not exclude near-duplicate positions or shared opening
structure.

\subsection{Position isolation and final-test firewall}

% receipt: evidence/benchmark_v2/exclusion_audit.complete.json
Candidate positions were rejected on exact placement and under the eight
dihedral board symmetries against 67 inventoried sources. The exclusion basis
contains 6,199,823 distinct placements and includes the 6,100,960-row Lichess
puzzle snapshot, prior render cohorts, training pools, and the named public
datasets. The benchmark and the separate adaptation pool are mutually
exact- and symmetry-disjoint. The audit treats material-only similarity as a
descriptive statistic, not as positional leakage. These checks concern canonical
piece placement against the supplied inventory. They do not establish coverage
of every public dataset, unknown foundation-model pretraining, or the complete
historical training lineage.

% receipt: evidence/benchmark_v2/benchmark.complete.json
The final-test partition has 100 position groups and 4,000 cell-images. Its
hashes and integrity metadata are frozen, but no recognition, localization,
attention, timing, or training result from that partition is reported here.
All model results in this paper use the development partition unless stated
otherwise. Here ``sealed'' names our no-scoring rule, not an access-control
or cryptographic barrier: releasing complete manifests, seeds, assets and a
deterministic generator would permit reconstruction of those images.

\subsection{Corner ground truth}

% receipt: evidence/benchmark_v2/benchmark.complete.json; validation scope
% source: tools/tests/test_paper_benchmark.py:360-447
The projected corners are produced by the same camera projection that renders
the scene. This avoids circular detector-derived ground truth. Validation
compares the analytic projection with renderer-recorded slab bounds and uses
asymmetric fixtures to test named-corner identity rather than only cyclic
self-consistency. The resulting corners support both perspective rectification
and detector-error measurements.

\subsection{Board conventions}

% receipt: evidence/docs/records/PARITY_CONVENTIONS.complete.json
The convention audit examines one representative image for each of the 56 image
sources and records named light corners together with loader behavior. The
project's earlier procedural lineage is consistently a1 light in the sampled
records, while the independently anchored public representatives for ChessReD,
ChessCog, CVChess, and ChessRender360 are a1 dark. The sampled SLCC frame has no
independent named-corner anchor, so its physical convention remains unresolved.
This is a source-sample audit, not a whole-dataset guarantee. None of the
reviewed loaders mirrors the position labels to compensate for board colour.
Appendix~\ref{app:conventions} gives the complete table.

\subsection{Metrics and aggregation}

% receipt: evidence/tools/reports/paper/lora_curve_2026-09-10/lora_report.json; metrics/dev
% receipt: evidence/tools/reports/paper/residual_shift_2026-09-10/r1/residual_r1.metrics.json
The primary benchmark results report exact-board and per-square class accuracy.
Exact-board accuracy requires every square label to match the target placement.
The separate release-checkpoint diagnostic also archives occupied-square class
accuracy, binary occupancy accuracy, confusion counts, and wrong-square
histograms; these are not measurements of the adapted checkpoints. Pooled
benchmark scores give equal weight to the registered cells, including the repeated
reference cells, rather than to distinct image hashes. Bootstrap comparisons
resample position groups with all their views together; the fixed synthetic
population and protocol still limit generalization.

% receipt: evidence/tools/reports/paper/FREEZE_2026-09-09.json; divergent_definition, splits
% source: tools/paper/annotate_benchmark.py:108-115; tools/paper/freeze_benchmark.py:203-243
For the earlier, unpaired benchmark, corner error is the mean Euclidean vertex
displacement after choosing the best cyclic alignment, divided by the mean
projected ground-truth edge length. Missing detections and vertices extending
more than two frame widths or heights outside the canvas are excluded from its
conditional summaries and counted separately. This normalization removes global
image scale, not projective foreshortening, and does not test named-corner
orientation. The earlier elevation plot is descriptive across different positions,
not a paired factor-effect estimate from the new benchmark.
